\chapter{Subdivisions}

This chapter follows Guibas--Stolfi, pp.~77--80, with enough intermediate
steps to read alongside the Lean proofs. The Lean files are the source of truth.
The ambient plane in Lean is $\mathbb C$, viewed as a real plane.
The local definition of a manifold does not include Hausdorffness; whenever a
result below needs that hypothesis, it is stated explicitly.

  % % % section for preliminaries leading up to definition 1.1
  \section{Preliminaries}

    % % % definition of separable, incident
    \begin{definition}[Separable, incident]\label{def:separable}
      \lean{GuibasStolfi.Separable, GuibasStolfi.Incident}
      \leanok
      Subsets $A$, $B$, of a topological space $M$ are \emph{separable} if some
      neighborhood of $A$ is disjoint from some neighborhood of $B$; otherwise
      they are \emph{incident}.
    \end{definition}

    \begin{lemma}\label{lem:incident-of-closure}
      \uses{def:separable}
      %\lean{GuibasStolfi.incident_of_mem_closure}

      If $A$ meets $\overline{B}$, then $A$ and $B$ are incident.
    \end{lemma}
    \begin{proof}

      Let $x \in A \cap \overline{B}$. Every neighbourhood of $A$ is a
      neighbourhood of $x$, so it meets every neighbourhood of $B$.
    \end{proof}

    % % % lemma relating paper's separable to mathlib4's separatedNhds
    \begin{lemma}\label{lem:separable-iff}
      \uses{def:separable}
      %\lean{GuibasStolfi.separable_iff_separatedNhds}

      $A$ and $B$ are separable if and only if they lie in disjoint open sets.
    \end{lemma}
    \begin{proof}

      Every neighborhood (open or closed) contains an open neighborhood
    \end{proof}

    % % % definitions of points, lines, disks
    \begin{definition}[Point, line, disk]\label{def:point-line-disk}
      \lean{GuibasStolfi.IsPoint, GuibasStolfi.IsLine, GuibasStolfi.IsDisk}
      \leanok
      A \emph{point} of $M$ is a one-element subset $\{ p \}$.
      A \emph{line} of $M$ is a subspace homeomorphic to $B_1 = (0, 1)$.
      A \emph{disk} of $M$ is a subspace homeomorphic to $B_2 = \{ x \in
      \mathbb{R}^2 : |x| < 1 \}$.
    \end{definition}

    % % % lemma distinguishing poitns, lines, disk
    \begin{lemma}\label{lem:shapes-not-homeomorphic}
      \uses{def:point-line-disk}
      \lean{GuibasStolfi.shapes_not_homeomorphic}
      \leanok
      No point is homeomorphic to a line or a disk, and no line is homeomorphic
      to a disk.
    \end{lemma}
    \begin{proof}
      A singleton $\{ p \}$ cannot be homeomorphic to a line or a disk, since
      each contains multiple distinct points.

      If a line and a disk were homeomorphic, then the open unit interval
      $(0, 1)$ and the plane would be homeomorphic. The removal of a point (say
      $\frac{1}{2}$) from the interval leaves the interval disconnected.
      Removing the image of the removed point in the plane under the
      homeomorphism (say $\varphi \left( \frac{1}{2} \right)$) leaves the plane
      connected. Such a homeomorphism cannot exist (since homeomorphisms
      preserve connectedness).
    \end{proof}

    % % % definition of 2 manifold
    \begin{definition}[2-manifold]\label{def:manifold}
      \uses{def:point-line-disk}
      \lean{GuibasStolfi.IsTwoManifold}
      \leanok
      A \emph{2-manifold} is a topological space in which every point has an
      open neighborhood homeomorphic to a disk.
    \end{definition}

    % % % lemma relating our 2 manifold to mathlib4's charted space
    \begin{lemma}\label{lem:manifold-charted}
      \uses{def:manifold}
      %\lean{GuibasStolfi.isTwoManifold_iff_chartedSpace}

      $M$ is a 2-manifold if and only if it admits a charted-space structure
      modeled on $\mathbb{R}^2$.
    \end{lemma}
    \begin{proof}

      An open disk is homeomorphic to $\mathbb{R}^2$. Conversely, every point
      of an open subset of $\mathbb{R}^2$ has neighbourhood that is a disk.
    \end{proof}

    \begin{definition}[Partition]\label{def:partition}
      %\lean{Setoid.IsPartition}

      \mathlibok
      A collection $C$ of subsets of $M$ is a \emph{partition} of $M$ if no
      member of $C$ is empty and every point of $M$ lies in exactly one member
      of $C$.
    \end{definition}

  % % % section on subdivisions
  \section{Subdivisions}

    % % % definition of simple path
    \begin{definition}[Simple path in $S^1$]\label{def:simple-path}
      \uses{def:partition}
      %\lean{GuibasStolfi.SimplePath}

      A \emph{simple path} in $S^1$ consists of distinct subsets
      $P_0,\ldots,P_m$, indexed cyclically modulo $m+1$, which partition $S^1$.
      Each $P_i$ is a singleton or an open subset of $S^1$ homeomorphic to
      $(0,1)$, and $P_i\cup P_{i+1}$ is connected for every $i$.
    \end{definition}
    Cyclic adjacency is explicit data in Lean. The word ``isolated'' refers
    to the point pieces of this partition, not to isolated points of the
    circle as a topological space. The one-point, one-arc partition is allowed.

    % % % definition 1.1, subdivision of a manifold
    \begin{definition}[Subdivision]\label{def:subdivision}
      \uses{def:manifold, def:point-line-disk, def:partition, def:simple-path}
      %\lean{GuibasStolfi.Subdivision}

      A \emph{subdivision} of a 2-manifold $M$ is a partition $S$ of $M$ into
      three finite sets of disjoint parts, the \emph{vertices} $\mathscr{V}
      \mathscr{S}$, the \emph{edges} $\mathscr{E} \mathscr{S}$, and
      \emph{faces} $\mathscr{F} \mathscr{S}$, with the following properties:
      \begin{description}
        % S1
        \item[S1] every vertex is a point of $M$

        % S2
        \item[S2] every edge is a line of $M$

        % S3
        \item[S3] every face is a disk of $M$

        % S4
        \item[S4] The boundary of every face is a closed path of edges and
                  vertices. Precisely, for every face $F$ there is a
                  simple path $\pi$ in $S^1$ and a continuous map
                  $\phi_F : \overline{B_2} \to \overline{F}$ onto
                  $\overline F$ that
                  \begin{description}
                    % % i % %
                    \item[(i)] $\phi_F$ is a homeomorphism from $B_2$ onto $F$

                    % % ii % %
                    \item[(ii)] maps each arc of $\pi$ homeomorphically onto an
                                edge of $S$

                    % % iii % %
                    \item[(iii)] maps each isolated point of $\pi$ to a vertex
                                 of $S$.
                  \end{description}
      \end{description}
      The vertices, edges, and faces of a subdivision are called its
      \emph{elements}.
    \end{definition}
    In Lean the map is defined on all of $\mathbb C$, with continuity required
    only on the closed disk; its values outside that disk are irrelevant.
    S4(ii) uses \emph{onto} an entire edge. This makes explicit the reading
    used by the subsequent proofs; the paper's literal wording is ``into''.
    Equivalence of these two formulations has not been proved here.

    % % % consequences of S4, paper pp. 77--78
      \begin{lemma}\label{lem:simple-path-alternate}
        \uses{def:simple-path}
        %\lean{GuibasStolfi.SimplePath.alternate}

        If $P_i$ is an arc of a simple path, then
        $\overline{P_i}=P_{i-1}\cup P_i\cup P_{i+1}$.
      \end{lemma}
      \begin{proof}

        Connectedness of consecutive unions forces each neighboring piece to
        be a singleton in the arc's closure: two disjoint nonempty open arcs
        would disconnect their union. A preconnected subset of the circle
        adds at most two points under closure, as can be seen by cutting the
        circle outside that subset and lifting it to an interval. Any further
        point of $\overline{P_i}$ would force the two neighboring singleton
        pieces to coincide. Injectivity of the indexed partition then reduces
        to the two-piece case, where no further piece exists. Thus these are
        exactly the neighbors added by closure. For one point and one arc,
        the two neighbors of the arc are the same singleton.
      \end{proof}

      \begin{lemma}\label{lem:arc-image}
        \uses{def:subdivision, lem:simple-path-alternate, def:separable}
        %\lean{GuibasStolfi.Subdivision.image_arc}

        If $\alpha$ is the arc between two consecutive points $a$ and $b$ of
        $\pi$, then its image $\phi_F(\alpha)$ is an edge incident to the
        points $\phi_F(a)$ and $\phi_F(b)$.
      \end{lemma}
      \begin{proof}

        \uses{lem:incident-of-closure}
        By (ii) and (iii), $\phi_F(\alpha)$ is an edge and $\phi_F(a)$,
        $\phi_F(b)$ are vertices. Since $a \in \overline{\alpha}$ and $\phi_F$
        is continuous, $\phi_F(a) \in \overline{\phi_F(\alpha)}$; apply
        \cref{lem:incident-of-closure}.
      \end{proof}

      \begin{definition}[Bounding path]\label{def:bounding-path}
        \uses{def:subdivision, lem:simple-path-alternate}
        %\lean{GuibasStolfi.Subdivision.boundingPath}

        The images of the elements of $\pi$, taken in the order in which they
        occur around $S^1$, form the \emph{bounding path} $\pi_F$ of $F$.
      \end{definition}

      \begin{lemma}\label{lem:face-boundary}
        \uses{def:bounding-path}
        %\lean{GuibasStolfi.Subdivision.frontier_face,
        %      GuibasStolfi.Subdivision.isConnected_frontier_face}

        $\pi_F$ is a closed, connected path of edges and vertices of $S$, whose
        union is the boundary $\overline{F} \setminus F$ of $F$.
      \end{lemma}
      \begin{proof}

        \uses{lem:arc-image}
        $\phi_F$ maps $\overline{B_2}$ onto $\overline{F}$ and $B_2$ onto $F$.
        The elements of $\pi$ map to vertices and edges, which are disjoint
        from $F$ since $S$ is a partition; so
        $\phi_F(S^1) = \overline{F} \setminus F$. It is connected as the
        continuous image of $S^1$ (Mathlib: \texttt{isConnected\_sphere},
        \texttt{IsConnected.image}).
      \end{proof}

      % Paper: pi_F need not be simple, since phi_F may take two or more
      % distinct arcs or points of
      % pi to the same element of S; so the closure of a face may not be a disk
      % (Fig. 1).
      Notice that the bounding path $\pi_F$ need not be simple, since $\phi_F$
      may take two or more distinct arcs or points of $\pi$ to the same element
      of $S$. Therefore the closure of a face may not be homeomorphic to a
      disk, as Figure~1 of the paper shows.

      \begin{lemma}\label{lem:no-finite-cover}
        \uses{def:point-line-disk, def:manifold}
        %\lean{GuibasStolfi.not_disk_subset_finite_union}

        In a Hausdorff topological space, a disk is not contained in a finite
        union of points and lines.
      \end{lemma}
      \begin{proof}

        Write each line as a countable union of compact interval images.
        Compact subsets are closed in the Hausdorff ambient space. Pull this
        countable cover back along a homeomorphism from the plane onto the
        disk. Baire's theorem on the plane gives a compact piece whose
        preimage contains a nonempty open ball. This produces a continuous
        injection of the plane into a point or into $(0,1)$. The former is
        impossible by cardinality. For the latter, removing a point of the
        plane leaves it connected, whereas an intermediate value between
        two image points cannot be omitted. This contradicts injectivity.
        No Baire assumption on the ambient space is needed.
      \end{proof}

      \begin{lemma}\label{lem:closures-cover}
        \uses{def:subdivision, def:separable}
        %\lean{GuibasStolfi.Subdivision.iUnion_closure_faces}

        Assume $M$ is Hausdorff. The closures of the faces cover $M$:
        $\bigcup_{F\in\FS}\overline F=M$.
      \end{lemma}
      \begin{proof}

        \uses{lem:no-finite-cover, lem:face-boundary, def:manifold, lem:incident-of-closure}
        $M \setminus \bigcup_F \overline{F}$ is open (finitely many faces) and
        covered by edges and vertices; if nonempty it contains a disk,
        contradicting \cref{lem:no-finite-cover}.
      \end{proof}

      \begin{lemma}\label{lem:edge-in-boundary}
        \uses{def:subdivision}
        %\lean{GuibasStolfi.Subdivision.edge_subset_frontier}

        Assume $M$ is Hausdorff. Every edge is entirely contained in the
        boundary of some face.
      \end{lemma}
      \begin{proof}

        \uses{lem:closures-cover, lem:face-boundary}
        Choose $x\in e$ and a face $F$ with $x\in\overline F$. The partition
        forces $x\notin F$. Lift $x$ to the boundary circle under $\phi_F$.
        Its path piece cannot map to a vertex, so it is an arc whose image
        is an edge containing $x$. Partition uniqueness identifies that
        edge with $e$. Surjectivity onto the whole edge gives the conclusion.
      \end{proof}

      \begin{definition}[Endpoints, loop]\label{def:endpoints}
        \uses{def:subdivision, def:separable}
        %\lean{GuibasStolfi.Subdivision.endpoints,
        %      GuibasStolfi.Subdivision.IsLoop}

        The \emph{endpoints} of an edge $e$ are the vertices of $S$ incident to
        $e$. If $e$ has only one endpoint, $e$ is a \emph{loop}.
      \end{definition}

      \begin{lemma}[Closure of an edge]\label{lem:closure-edge}
        \uses{lem:edge-in-boundary, lem:simple-path-alternate, def:subdivision}
        %\lean{GuibasStolfi.Subdivision.closure_edge}

        Assume $M$ is Hausdorff. For every edge $e$ there are vertices
        $\{p\},\{q\}$ such that $\overline e=e\cup\{p,q\}$.
      \end{lemma}
      \begin{proof}

        Write $e=\phi_F(P_i)$ for a boundary arc. Its closure in the circle
        is the arc together with its two neighboring singleton pieces.
        Continuity on this compact closure and Hausdorffness imply
        $\overline{\phi_F(P_i)}=\phi_F(\overline{P_i})$.
        S4 maps the two singleton pieces to vertices; their images may agree.
      \end{proof}

      \begin{lemma}\label{lem:endpoints-card}
        \uses{def:endpoints}
        %\lean{GuibasStolfi.Subdivision.endpoints_card}

        Assume $M$ is Hausdorff. Every edge is incident to two (not
        necessarily distinct) vertices of $S$.
      \end{lemma}
      \begin{proof}

        \uses{lem:closure-edge, lem:compact, lem:arc-image}
        By \cref{lem:edge-in-boundary} $e = \phi_F(\alpha)$ for an arc $\alpha$
        between points $a, b$, and $\phi_F(a)$, $\phi_F(b)$ are endpoints
        (\cref{lem:arc-image}). Conversely, $\overline{e} =
        \phi_F(\overline{\alpha}) = e \cup \{\phi_F(a), \phi_F(b)\}$ (compact
        image in a Hausdorff space), and a vertex outside $\overline{e}$ is
        separable from $e$ by regularity.
      \end{proof}

      \begin{lemma}\label{lem:loop-circle}
        \uses{def:subdivision, def:endpoints}
        %\lean{GuibasStolfi.Subdivision.closure_loop_homeomorph_circle}

        Assume $M$ is Hausdorff. The closure of a loop is homeomorphic to $S^1$.
      \end{lemma}
      \begin{proof}

        \uses{lem:endpoints-card, lem:closure-edge, lem:compact}
        The two vertices added under closure must be the same unique
        endpoint, so $\overline e=e\cup\{p\}$ with $p\notin e$.
        This is a compact Hausdorff space containing $e\cong(0,1)$ as the
        complement of one point. It is therefore the one-point compactification
        of $(0,1)$, homeomorphic to $S^1$. This also covers the one-point,
        one-arc boundary path without treating its closure as a closed interval.
      \end{proof}

      \begin{lemma}\label{lem:compact}
        \uses{def:subdivision}
        %\lean{GuibasStolfi.Subdivision.compactSpace}

        Assume $M$ is Hausdorff. Then $M$ is compact.
      \end{lemma}
      \begin{proof}

        \uses{lem:closures-cover}
        By \cref{lem:closures-cover}, the face closures cover $M$, and $\overline{F} =
        \phi_F(\overline{B_2})$ is compact. Thus $M$ is a finite union of
        compact sets. The Mathlib ingredients are:
        \begin{quote}\small
          \texttt{isCompact\_closedBall}, \texttt{IsCompact.image},\\
          \texttt{Set.Finite.isCompact\_biUnion}.
        \end{quote}
      \end{proof}

  % % % section 1.1 in paper, Equivalence and Connectivity
  \section{Equivalence and Connectivity}

    The isomorphism definitions and incidence invariance below do not require
    Hausdorffness. The graph and connectivity results explicitly use it.

    % % % definition 1.2, subdivision isomorphisms and equivalent subdivisions
    \begin{definition}[Isomorphism, equivalent]\label{def:isomorphism-equivalent}
      \uses{def:subdivision}
      %\lean{GuibasStolfi.Subdivision.Iso, GuibasStolfi.Subdivision.Equivalent}

      Let $S$ and $S'$ be two subdivisions of the manifolds $M$ and $M'$. An
      \emph{isomorphism} from $S$ to $S'$ is a homeomorphism of $M$ onto $M'$
      that maps each element of $S$ onto an element of $S'$. When such a
      mapping exists, we say that $S$ and $S'$ are \emph{equivalent}, and we
      write $S \sim S'$.
    \end{definition}

    \begin{lemma}\label{lem:shapes-distinct}
      \uses{def:point-line-disk}
      %\lean{GuibasStolfi.not_homeomorphic_point_line,
      %      GuibasStolfi.not_homeomorphic_point_disk,
      %      GuibasStolfi.not_homeomorphic_line_disk}

      No point is homeomorphic to a line or a disk, and no line is homeomorphic
      to a disk.
    \end{lemma}
    \begin{proof}

      A point has one element; lines and disks are infinite. Removing any point
      disconnects a line, but no disk: $B_2 \cong \mathbb{R}^2$, and
      $\mathbb{R}^2 \setminus \{x\}$ is connected (Mathlib:
      \texttt{isConnected\_compl\_singleton\_of\_one\_lt\_rank}).
    \end{proof}

    \begin{lemma}\label{lem:iso-preserves-type}
      \uses{def:isomorphism-equivalent}
      %\lean{GuibasStolfi.Subdivision.Iso.mapsTo_vertices,
      %      GuibasStolfi.Subdivision.Iso.mapsTo_edges,
      %      GuibasStolfi.Subdivision.Iso.mapsTo_faces,
      %      GuibasStolfi.Subdivision.Iso.symm}

      An isomorphism maps vertices onto vertices, edges onto edges, and faces
      onto faces.
    \end{lemma}
    \begin{proof}

      \uses{lem:shapes-distinct}
      The image of an element is an element of $S'$ homeomorphic to it; apply
      \cref{lem:shapes-distinct}. The inverse homeomorphism also maps elements
      to elements: pull back one point of a target element and use partition
      uniqueness to identify its whole preimage cell. Applying the same
      shape argument to the inverse gives the bijections of cell families.
    \end{proof}

    \begin{lemma}\label{lem:iso-preserves-incidence}
      \uses{def:isomorphism-equivalent, def:separable}
      %\lean{GuibasStolfi.Subdivision.Iso.incident_iff}

      An isomorphism preserves incidence: $A, B$ are incident iff $h(A), h(B)$
      are.
    \end{lemma}
    \begin{proof}

      A homeomorphism maps neighbourhoods to neighbourhoods, so it preserves
      separability.
    \end{proof}

    % % % definition of......... topological property of subdivisions
    \begin{definition}[Topological property]\label{def:topological-property}
      \uses{def:subdivision, def:isomorphism-equivalent}
      %\lean{GuibasStolfi.Subdivision.TopologicalProperty}

      A \emph{topological property} of subdivisions is a property that is
      invariant under equivalence.
    \end{definition}

    % % % definition for graph of a subdivision
    \begin{definition}[Graph of subdivision]\label{def:graph}
      \uses{def:subdivision, def:endpoints}
      %\lean{GuibasStolfi.Subdivision.graph}

      Assume $M$ is Hausdorff. The graph of $S$ has vertex set $\VS$ and
      edge set $\ES$; an edge $e$ joins $u,v$ exactly when its endpoint set is
      $\{u,v\}$. Thus loops and distinct edges with the same endpoints are allowed.
    \end{definition}

    % Stand-in: Mathlib's Graph has no connectivity yet
    % (only Connected/EdgeCut.lean).
    \begin{definition}[Connected graph]\label{def:graph-connected}
      %\lean{GuibasStolfi.Graph.Connected}

      A graph is \emph{connected} if any two vertices $u, v$ are joined by a
      finite sequence of vertices $u = v_0, v_1, \dots, v_n = v$ in which
      $v_{i-1}$ and $v_i$ are the ends of some edge.
    \end{definition}

    % Stand-in: Mathlib's Graph has no morphisms yet
    % (Combinatorics/Graph/Maps.lean, TODO).
    \begin{definition}[Graph isomorphism]\label{def:graph-iso}
      %\lean{GuibasStolfi.Graph.Iso}

      An \emph{isomorphism} of graphs is a pair of bijections, one between the
      vertex sets and one between the edge sets, such that an edge joins $u$
      and $v$ if and only if its image joins the images of $u$ and $v$.
    \end{definition}

    % % % lemma that equivalent subdivisions have isomorphic graphs
    \begin{lemma}\label{lem:equiv-graph-iso}
      \uses{def:subdivision, def:isomorphism-equivalent, def:graph, def:graph-iso}
      %\lean{GuibasStolfi.Subdivision.Equivalent.graphIso}

      For Hausdorff ambient spaces, the graphs of two equivalent subdivisions
      $S$ and $S'$ are isomorphic.
    \end{lemma}
    \begin{proof}

      \uses{lem:iso-preserves-type, def:endpoints, lem:iso-preserves-incidence}
      Let $h$ be an isomorphism. By \cref{lem:iso-preserves-type},
      $v \mapsto h(v)$ and $e \mapsto h(e)$ are bijections
      $\mathscr{V}\mathscr{S} \to \mathscr{V}\mathscr{S}'$ and
      $\mathscr{E}\mathscr{S} \to \mathscr{E}\mathscr{S}'$. Endpoints are
      preserved because they are the incident vertices of the edge, and the
      homeomorphism preserves incidence in both directions.
    \end{proof}

    % % % COUNTEREXAMPLE FOR CONVERSE: SEE FIGURE 2 IN PAPER, PG. 79
    \begin{proposition}[Converse fails]\label{prop:graph-not-determining}
      \notready
      \uses{prop:torus-subdivision, prop:klein-subdivision, prop:torus-klein-graph,
            prop:torus-klein-nonequivalent}
      There exist subdivisions $S$, $S'$ whose graphs are isomorphic but
      $S \not\sim S'$ (Fig.~2).
    \end{proposition}
    The two subdivision instances and their graph isomorphism are now
    formalized later in this chapter. Nonequivalence remains open; the proposition
    is not marked proved merely because its examples and graph map exist.

    \begin{lemma}\label{lem:faces-open}
      \uses{def:subdivision, def:separable}
      %\lean{GuibasStolfi.Subdivision.isOpen_face}

      Assume $M$ is Hausdorff. Every face is open in $M$; in particular
      distinct faces are separable.
    \end{lemma}
    \begin{proof}

      \uses{lem:face-boundary, lem:closures-cover, lem:endpoints-card, lem:separable-iff}
      Let $F$ be a face. Every point of $M$ lies in exactly one element of $S$,
      so
        \[
          M \setminus F \;=\; \bigcup_{v \in \mathscr{V}\mathscr{S}} v
            \;\cup\; \bigcup_{e \in \mathscr{E}\mathscr{S}} \overline{e}
            \;\cup\; \bigcup_{F' \in \mathscr{F}\mathscr{S},\, F' \neq F}
            \overline{F'},
        \]
        provided each set on the right is disjoint from $F$. Vertices are
        disjoint from $F$ since $S$ is a partition. For an edge $e$,
        $\overline{e}$ consists of $e$ and its endpoints
        (\cref{lem:endpoints-card}), all elements other than $F$. For a face
        $F' \neq F$, $\overline{F'}$ is $F'$ together with the elements of
        $\pi_{F'}$ (\cref{lem:face-boundary}), again all elements other than
        $F$. Each set on the right is closed: vertices are points, which are
        closed since $M$ is Hausdorff, and closures are closed. There are
        finitely many, so $M \setminus F$ is closed and $F$ is open. Distinct
        faces are then disjoint open sets, hence separable
        (\cref{lem:separable-iff}).
    \end{proof}

    \begin{lemma}\label{lem:face-incident-edge}
      \uses{def:subdivision, def:separable}
      %\lean{GuibasStolfi.Subdivision.exists_edge_incident_face}

      Every face is incident to some edge.
    \end{lemma}
    \begin{proof}

      \uses{lem:incident-of-closure}
      $S^1$ is not a finite set of points, so $\pi$ has an arc; its image is an
      edge contained in $\overline{F}$ by (ii); apply
      \cref{lem:incident-of-closure}.
    \end{proof}

    \begin{lemma}\label{lem:incident-union-connected}
      \uses{def:separable}
      %\lean{GuibasStolfi.IsConnected.union_of_incident}

      In a completely normal space, if $A$ and $B$ are connected and incident,
      then $A \cup B$ is connected.
    \end{lemma}
    \begin{proof}

      If $x\in A\cap\overline B$, adjoining $x$ to $B$ preserves
      preconnectedness, and the enlarged set meets $A$. Their union is
      connected. The case $x\in\overline A\cap B$ is symmetric.
      Otherwise $A,B$ are separated sets, and complete normality supplies
      disjoint open neighborhoods, contradicting incidence.
    \end{proof}

    \begin{lemma}[Complete normality of a subdivided surface]\label{lem:completely-normal}
      \uses{lem:manifold-charted, lem:compact}
      %\lean{GuibasStolfi.Subdivision.completelyNormalSpace}

      A subdivision with Hausdorff ambient space has completely normal
      ambient space.
    \end{lemma}
    \begin{proof}

      Its local disk neighborhoods supply a charted-space structure.
      A compact Hausdorff manifold modeled on the plane is metrizable;
      metrizability implies complete normality.
    \end{proof}

    % % % Theorem 1.1 in paper
    \begin{theorem}\label{thm:graph-connected-iff}
      \uses{def:graph, def:graph-connected}
      %\lean{GuibasStolfi.Subdivision.graph_connected_iff}

      Assume $M$ is Hausdorff. The graph of $S$ is connected if and only if
      $M$ is preconnected, that is, $M$ has no separation into two nonempty
      disjoint open sets. Both sides allow the empty case. For nonempty $M$
      this agrees with the usual connected-space convention.
    \end{theorem}
    \begin{proof}

      \uses{lem:faces-open, lem:face-boundary, lem:face-incident-edge,
            lem:incident-union-connected, lem:completely-normal}
      If the graph is connected, attach to each point a connected union of
      cells containing a vertex. Join the resulting vertices by a graph path;
      the union of these connected sets contains the original pair of points.
      Incidence joins connected sets using complete normality.

      Conversely, suppose two vertices are not joined. Group the vertices
      and edges incident to vertices reachable from the first vertex, then
      add all faces incident to that group. The group and its complement are
      separable: incident vertices or edges have a common incident vertex;
      a face cannot meet both groups on its connected boundary; distinct
      faces are separable. Finiteness combines the pairwise separations into
      disjoint open neighborhoods of the two groups. These cover $M$ and
      each contains one of the chosen vertices, contradicting preconnectedness.
    \end{proof}

    \begin{lemma}\label{lem:restrict-component}
      \uses{def:subdivision}
      %\lean{GuibasStolfi.Subdivision.restrictComponent}

      Assume $M$ is Hausdorff. Let $C$ be a connected component of $M$.
      The elements of $S$ contained
      in $C$ form a subdivision of $C$.
    \end{lemma}
    \begin{proof}

      \uses{def:manifold}
      Every element of $S$ is connected (a point, line or disk), so it lies in
      exactly one component; hence the elements inside $C$ partition $C$. $M$
      is locally connected, so $C$ is open (Mathlib:
      \texttt{isOpen\_connectedComponent}) and therefore a 2-manifold.
      (S1)--(S3) are unchanged, and for a face $F \subseteq C$, $\overline{F}$
      is connected and meets $C$, so $\phi_F$ lands in $C$ and (S4) holds.
    \end{proof}

    \begin{corollary}\label{cor:components}
      \uses{def:graph-connected, def:graph}
      %\lean{GuibasStolfi.Subdivision.componentsEquiv}

      Assume $M$ is Hausdorff. The connected components of the manifold are in one-to-one
      correspondence with the connected components of the underlying graph.
    \end{corollary}
    \begin{proof}

      \uses{thm:graph-connected-iff, lem:restrict-component}
      Each component $C$ carries the subdivision $S|_C$
      (\cref{lem:restrict-component}), whose graph is the part of the graph of
      $S$ inside $C$, and is connected by \cref{thm:graph-connected-iff}. An
      edge lies in one component together with its endpoints, so no edge joins
      the graphs of different components. Hence $C \mapsto$ (graph of $S|_C$)
      is a bijection onto the components of the graph.
    \end{proof}
